\documentclass[11pt]{amsart}

\usepackage{amsmath,amssymb,amsthm}
\usepackage{hyperref}

\newtheorem{theorem}{Theorem}[section]
\newtheorem{lemma}[theorem]{Lemma}

\newcommand{\hs}{\boxplus}
\newcommand{\cP}{\mathcal P}

\title{A Cubic Consequence of a Vanishing Ternary Hypersum}
\date{}

\begin{document}

\maketitle

\section{Conventions and preliminary lemmas}

We use the standard commutative notion of a hyperfield. Thus
\((\mathbb H,\hs,0)\) is a commutative, associative, reversible hypergroup
with identity \(0\), the nonzero elements form an abelian group under
multiplication, multiplication by \(0\) is \(0\), and multiplication
distributes over hyperaddition:
\[
a(b\hs c)=ab\hs ac.
\]
The proof below uses only the corresponding inclusion
\[
a(b\hs c)\subseteq ab\hs ac,
\]
so it also applies under the weaker distributivity convention.

For subsets \(A,B\subseteq\mathbb H\), write
\[
A\hs B=\bigcup_{\substack{a\in A\\ b\in B}}a\hs b.
\]
Associativity and commutativity make finite hypersums such as
\(a_1\hs\cdots\hs a_n\) independent of parentheses and ordering.

\begin{lemma}[Additive inverses]
\label{lem:additive-inverses}
For every \(a,b\in\mathbb H\),
\[
0\in a\hs b
\quad\Longleftrightarrow\quad
b=-a.
\]
Moreover,
\[
-a=(-1)a.
\]
In particular, \((-a)^3=-(a^3)\).
\end{lemma}

\begin{proof}
If \(0\in a\hs b\), reversibility gives
\[
b\in 0\hs(-a)=\{-a\},
\]
so \(b=-a\). Conversely, if \(b=-a\), then \(0\in a\hs(-a)=a\hs b\)
by the inverse axiom.

Next, \(0\in1\hs(-1)\). Multiplication by \(a\), followed by
distributivity, gives
\[
0\in a\bigl(1\hs(-1)\bigr)\subseteq a\hs(-1)a.
\]
Hence the first part yields \((-1)a=-a\). Finally,
\[
(-a)^3=\bigl((-1)a\bigr)^3=(-1)^3a^3=(-1)a^3=-(a^3).
\qedhere
\]
\end{proof}

\begin{lemma}[Cubic expansion for an element of a binary hypersum]
\label{lem:cubic-expansion}
Let \(x,y,s\in\mathbb H\), and suppose that \(s\in x\hs y\). Put
\[
T=x^2y\hs xy^2.
\]
Then
\[
s^3\in x^3\hs y^3\hs T\hs T\hs T.
\]
\end{lemma}

\begin{proof}
Since \(s\in x\hs y\), distributivity gives
\[
sx\in x^2\hs xy,
\qquad
sy\in xy\hs y^2.
\]
Also,
\[
s^2\in s(x\hs y)\subseteq sx\hs sy.
\]
Consequently,
\[
s^2\in (x^2\hs xy)\hs(xy\hs y^2).
\]
Multiplying by \(s\) and again applying distributivity yields
\[
s^3
\in
(sx^2\hs sxy)\hs(sxy\hs sy^2).
\]
Using \(s\in x\hs y\) once more, we obtain
\[
sx^2\in x^3\hs x^2y,\qquad
sxy\in x^2y\hs xy^2,\qquad
sy^2\in xy^2\hs y^3.
\]
Therefore
\[
s^3
\in
(x^3\hs x^2y)
\hs
(x^2y\hs xy^2)
\hs
(x^2y\hs xy^2)
\hs
(xy^2\hs y^3).
\]
By associativity and commutativity of hypersaddition, the latter set is
\[
x^3\hs y^3
\hs
(x^2y\hs xy^2)
\hs
(x^2y\hs xy^2)
\hs
(x^2y\hs xy^2),
\]
which is precisely \(x^3\hs y^3\hs T\hs T\hs T\).
\end{proof}

\begin{lemma}[The mixed quadratic term]
\label{lem:mixed-term}
Suppose that
\[
0\in x\hs y\hs z.
\]
Then
\[
-xyz\in x^2y\hs xy^2.
\]
\end{lemma}

\begin{proof}
Multiplying the given hypersum by \(xy\) and using distributivity gives
\[
0\in xy(x\hs y\hs z)
\subseteq x^2y\hs xy^2\hs xyz.
\]
Thus there exists
\[
t\in x^2y\hs xy^2
\]
such that \(0\in t\hs xyz\). By Lemma~\ref{lem:additive-inverses},
\(t=-xyz\). Hence
\[
-xyz\in x^2y\hs xy^2.
\qedhere
\]
\end{proof}

\section{Main theorem}

\begin{theorem}
\label{thm:main-cubic}
Let \(\mathbb H\) be a hyperfield, and let \(x,y,z\in\mathbb H\). If
\[
0\in x\hs y\hs z,
\]
then
\[
0\in
x^3\hs y^3\hs z^3
\hs(-xyz)\hs(-xyz)\hs(-xyz).
\]
\end{theorem}

\begin{proof}
By associativity, there exists an element \(s\in x\hs y\) such that
\[
0\in s\hs z.
\]
Lemma~\ref{lem:additive-inverses} gives \(s=-z\). Hence
\[
-z\in x\hs y.
\]

Let
\[
T=x^2y\hs xy^2
\]
and
\[
C=x^3\hs y^3\hs T\hs T\hs T.
\]
Lemma~\ref{lem:cubic-expansion}, applied to \(s=-z\), gives
\[
(-z)^3\in C.
\]
By Lemma~\ref{lem:additive-inverses}, \((-z)^3=-z^3\), so
\[
-z^3\in C.
\]

Put
\[
t=-xyz.
\]
Lemma~\ref{lem:mixed-term} gives \(t\in T\). Therefore
\[
x^3\hs y^3\hs t\hs t\hs t
\subseteq
x^3\hs y^3\hs T\hs T\hs T
=C.
\]
Since \(-z^3\in C\), it follows that
\[
-z^3\in x^3\hs y^3\hs t\hs t\hs t.
\]
Consequently,
\[
0\in(-z^3)\hs z^3
\subseteq
x^3\hs y^3\hs z^3\hs t\hs t\hs t.
\]
Substituting \(t=-xyz\) yields
\[
0\in
x^3\hs y^3\hs z^3
\hs(-xyz)\hs(-xyz)\hs(-xyz).
\qedhere
\]
\end{proof}

Thus the answer to the question is affirmative.

\end{document}